\chapter{Results}
\label{ch:results}

\keybox{how to read this chapter}{%
The chapter runs in the order the work was done, which is also the order in which
each claim depends on the one before it.

First the reproduction: can this project produce the two published baselines at all? Nothing later means
anything without it. Then the two instruments, which measure the capture geometry itself and never a rendered image --- these are the
falsification gates, and they were run before the method existed. Then the method
on the standard benchmark, where the theory demands it change nothing. Then the
stress suite, where the theory says it should act, which is the decisive
experiment. Then the second method, which fails.

Three conventions hold throughout. Every number here is read from a single
append-only record and inserted mechanically; nothing is typed by hand,
and a quantity that has not been measured appears as a visible marker, never as a plausible figure. Published figures are cited as \emph{targets} and never
used as a control --- every comparison is against this project's own re-run
baseline. And where something went wrong, it is reported where it happened,
with the correction it forced; the engineering detail behind each is in
Appendix~\ref{app:defects} so that it does not interrupt the argument.}

\keybox{status}{%
Five reproduction rungs, two read-only instruments and both proposed methods
have been run, and everything below is measured. The protocol, acceptance
criteria, compute budget and decision gates were frozen \emph{before} any run,
so that a disappointing number later could not be rationalised into a success;
where a gate failed or a prediction did not hold, the chapter says so and keeps
the original wording. Every table and figure in this chapter is regenerated
mechanically from the measurement record; no number is typed into this document
by hand (Appendix~\ref{app:record}).}

\section{Pre-flight assertions}
\label{sec:preflight}

A reproduction that has not checked its own instrument is a reproduction of
nothing in particular. Twelve assertions therefore run before the first training
step, and the session aborts if any of them fails. They check that the GPU is
one the code can actually run on, that both arms load the same images, that the
perceptual metric uses the backbone the papers used, and that the two checkouts differ by exactly the intended change and
nothing else. Appendix~\ref{app:harness} lists all twelve with what each one
protects against. All twelve pass.
\section{Single-scale training}
\label{sec:table1}

\input{generated/table1-stmt}

\keybox{what reproduces, and to what level}{%
Over the $\stmtArmAScenes$ scenes the published mean covers, Mip-Splatting lands
within $\stmtArmAMaxDelta$\,dB of its published figure at every test scale.
That is inside the $\pm0.20$\,dB of Table~\ref{tab:levels}, so the Mip-Splatting arm carries an
\textbf{L1} claim on PSNR. The \textbf{gap} between the arms, on the
$\stmtMatchedScenes$ scenes both measured, is within $\stmtGapMaxDelta$\,dB of
the published gap at all four scales --- also L1, and it is the quantity the
two-arm construction exists to produce.

The 3DGS column is a different matter and claims nothing. \texttt{ship} does not
train on the vanilla arm (\S\ref{sec:shipfails}), so that column is a
$\stmtArmBScenes$-scene mean standing against an eight-scene published one. Its $\stmtArmBMaxDelta$\,dB is mostly the missing scene, and is reported as
such.}

\keybox{an artefact in an earlier version of this table}{%
Both \emph{ours} columns were computed on the matched seven scenes, and compared
against published figures that are means over eight. \texttt{ship} happens to be
the Mip-Splatting arm's weakest scene at full resolution, at $30.48$\,dB against a set mean of
$33.44$, so excluding it lifted the Mip-Splatting arm's mean by $0.42$\,dB, and that entire
lift was being reported as reproduction error. The Mip-Splatting arm read $0.50$\,dB from its
published figure and was written up as L2.

Two comparisons live in this table, and each needs its own denominator. A
column compared against a published eight-scene mean has to be an eight-scene
mean. A gap between two arms has to be over the scenes both arms have, or it is
partly the difference between two sets of objects. Computing both on the
matched set made the second correct and the first wrong, and the error was
invisible because it moved the number in the direction a reproduction error
would move it.}

\subsection*{The same result, as a picture}
\label{sec:qualitative}

A decibel is a summary, and a $10$\,dB gap in a table is a claim the reader has
to take on trust. Figure~\ref{fig:ladder} is the first row of
Table~\ref{tab:stmt} rendered instead of tabulated: one view of \texttt{lego},
trained once at full resolution, evaluated at each of the four test scales, on
both arms.

\begin{figure}[htbp]
  \centering
  \includegraphics[width=0.92\linewidth]{figures/fig11-qualitative-ladder.pdf}
  \caption{One view of \texttt{lego}, single-scale training, four test scales.
  Left: ground truth at that scale. Middle: 3DGS. Right: Mip-Splatting. The
  numbers are PSNR for the two images shown, recomputed from those pixels, which is why they differ from Table~\ref{tab:stmt}, a mean
  over $\stmtMatchedScenes$ scenes and every test view. Renders are magnified
  with nearest-neighbour interpolation so that a rendered pixel is a visible
  block: a $\tfrac18$-scale render is $100\times100$, and any smooth resampler
  would average away the aliasing the figure exists to show.}
  \label{fig:ladder}
\end{figure}

At the training scale the two are indistinguishable, and the numbers agree to
$0.07$\,dB. One octave down they have separated by $7$\,dB, and what separates
them is visible without being told where to look: the 3DGS reconstruction
\emph{thickens}. Treads merge into a band, the gap between the shovel arms
closes, the yellow spreads past the studs that bound it. By $\tfrac18$ the
machine is a blob with a red dot on it, and Mip-Splatting still resolves the
tracks.

That thickening is the mechanism of Chapter~\ref{ch:background} in visible
form. A primitive smaller than a pixel at the test scale is still splatted, at full
opacity, into the pixel that contains it, and many such primitives sum. The image gains energy it should not have, so
edges bloom outward and dark regions lift. The band-limit is what stops a
primitive from being smaller than the sampling can support, and the right-hand
column is what that buys.

\begin{figure}[htbp]
  \centering
  \includegraphics[width=0.92\linewidth]{figures/fig12-qualitative-scenes.pdf}
  \caption{The $\tfrac18$ row across four scenes, same construction as
  Figure~\ref{fig:ladder}. \texttt{materials} is the clearest case: the
  spheres lose their specular highlights to bloom and the dark ones lift
  towards grey. Every scene measured under both arms shows the same signature;
  these four are shown because they are visually distinct from one another. The
  remaining scenes are kept as per-scene contact sheets alongside the record, so that the selection can be audited instead of trusted.}
  \label{fig:qualscenes}
\end{figure}

\keybox{what these figures show}{%
They are the pixels the PSNR column was computed from. Each rung deletes all but
the first few rendered views once the metrics have consumed them, and what
survives is a matched set of predictions and ground truths per model; these are those files, carried straight out of the run that produced them. The ground-truth image is
byte-identical between the two arms at every index, and that is checked on every
pair, so the two predictions are of the same view.

They are also \emph{single views}, and a single view is not evidence of a mean.
The tables remain the claim; these show what the claim looks like.}

\input{generated/table-cost}

\section{The scale-collapse gate}
\label{sec:g2}

The scale-collapse gate was declared in advance as: 3DGS falls $\gtrsim13$\,dB from full to
$\tfrac18$; Mip-Splatting falls $\lesssim6$\,dB; the gap at $\tfrac18$ exceeds
$9$\,dB. On R1 it failed, on two of its three criteria: the 3DGS arm fell $9.09$\,dB
where $13$ was required, and the gap at $\tfrac18$ was $4.07$\,dB against a
published $10.98$. The protocol's instruction when the scale-collapse gate fails is to stop and
diagnose in a stated order. That took two attempts, and the first was wrong.

\subsection*{The protocol's three checks, and a fourth}

\begin{enumerate}[leftmargin=1.4em,itemsep=2pt]
  \item \textbf{The per-scale recovery.} Its count-weighted mean reproduces
        the shipped evaluation's own pooled figure to within $0.01$\,dB on every
        scene of every run --- asserted per scene, on every scene. Not the
        cause.
  \item \textbf{The 3DGS arm's \texttt{filter\_3D} at step 30\,000.} Read back from all
        eight saved point clouds it is identically zero, and non-zero on all
        eight of the Mip-Splatting arm's. The reduction is exact. Not the cause.
  \item \textbf{The white-background setting.} Passed to every job on both
        arms. Not the cause.
\end{enumerate}

\paragraph{The first diagnosis, and why it was wrong.} With those three
eliminated, the explanation offered was densification. The public Mip-Splatting repository
changed how it decides when to subdivide a primitive, roughly a year after the
paper was published: it adopted a criterion from Gaussian Opacity Fields, which
splits on a quantile over accumulated gradient magnitudes instead of on a fixed
gradient threshold. The original
3DGS rule is the fixed threshold, $\lVert\nabla\rVert \ge \tau_{\mathrm{pos}}$
alone. The argument was that a
better densification recovers most of what the missing band-limit costs, so the collapse the gate's thresholds were calibrated against could not occur in
this codebase, and the gate was therefore unmeetable here.

It was a plausible story, it was consistent with every number then available,
and it was \emph{wrong}. It also had the shape a wrong diagnosis usually has:
it explained the failure as unfixable, and it did so without a test that could
have refuted it.

\subsection*{The fourth check: the rasteriser on the 3DGS arm}
\label{sec:c1}

The three checks above say the gate failed. They do not say why, and the first
answer this project gave was wrong.

The 3DGS arm was built by zeroing Mip-Splatting's 3D smoothing filter, which makes the
world-space band-limit 3DGS's. What that overlooks is that Mip-Splatting has a
\emph{second} band-limit: a screen-space dilation carrying an opacity
compensation $\rho = \sqrt{\det\Sigma' / \det(\Sigma' + k I)}$. That one lives in
CUDA,
where zeroing a Python attribute does not reach it. The arm labelled 3DGS was therefore anti-aliasing better than 3DGS, at the very scales the gate was
scored on.

Isolating that term is what the opacity-compensation fix does. Appendix~\ref{app:impl} gives the switch and
what it touches; Appendix~\ref{app:defects} gives how it was found, including the
two further things the switch turned out to remove that nobody had asked it to.
\subsection*{The scale-collapse gate, re-scored}

Every 3DGS-arm row produced with the compensation active is superseded. They are
marked as such in the record and retained --- every table in this document
excludes them, and none of them has been deleted.
The Mip-Splatting arm is untouched and was not re-run: nothing about this defect involves it.

\input{generated/table-g2}

All three criteria pass, against thresholds that were set before any run and have
not been touched since. And the consequence reaches further than the gate: with
the arm rebuilt, Table~\ref{tab:stmt} reproduces \emph{both} baselines, where before it
reproduced one. 3DGS lands within $\stmtArmBMaxDelta$\,dB of its published figure at
every test scale and Mip-Splatting within $\stmtArmAMaxDelta$, and the gap column
--- the quantity the whole two-arm construction exists to measure --- reproduces
to within $\stmtGapMaxDelta$\,dB at all four, which is inside the L1 tolerance. This project did not have that before; it had the Mip-Splatting arm at L1 and a 3DGS arm several decibels adrift, for a
reason it had explained incorrectly.

\keybox{the cost of the rebuild}{%
It cost the R1 and R2 3DGS-arm columns and the 3DGS half of the seed spread --- three rungs of measurement, re-run, and it surfaced a numerical
fragility that had been hidden by the compensation (\S\ref{sec:shipfails}). It bought a 3DGS arm that is 3DGS in both of its band-limits, so the control every later comparison rests on is the method itself. The
between-arm comparison was the one thing the earlier diagnosis insisted was
safe, and it was the thing that was wrong.}

\subsection*{What the vanilla arm cost in robustness}
\label{sec:shipfails}

Removing the compensation cost something the protocol had not budgeted for. The
scene \texttt{ship} is much the largest in the suite, at $442\,235$ primitives
and $90\,\%$ of the card, and it does not train on the vanilla arm at all. It
has
been attempted nine times and has failed nine times, always with a CUDA illegal
memory access and never with an allocation error, at iterations spread from
$3\,000$ to $12\,100$. What varies is \emph{when}, not \emph{whether}.

Two things keep the exclusion narrow. \texttt{ship} trains
without incident on the \emph{multi-scale} protocol, so Table~\ref{tab:mtmt} is a
full eight-scene comparison and the failure is specific to single-scale training.
And it trains without incident on the Mip-Splatting arm, so the fragility belongs to the vanilla band-limit semantics.

It is excluded from the single-scale 3DGS-arm rows, and the protocol's rule for
a scene that will not run is to record it and move on. The instrumentation that
localised the fault, and the three explanations it ruled out, are in
Appendix~\ref{app:defects}.
\section{Multi-scale training}
\label{sec:table2}

\input{generated/table2-mtmt}

\paragraph{This column was two rasterisers, briefly.} The 3DGS arm was re-measured
after the opacity-compensation fix, and for a period each of its $32$ cells was the mean of the old measurement and the new one instead of the new one alone --- the rows the re-run
replaced were never retired, because the rule that retires them keyed on a name
the re-run shared. Nothing about the output looked wrong. The generator now
refuses to emit a table whose cells draw on more than one build, and names the
cell; the account is in Appendix~\ref{app:defects}.
\section{The densification explanation}
\label{sec:densification}

Section~\ref{sec:g2} settles the cause of the $\tfrac18$-scale shortfall: it was
the 2D opacity compensation, and removing it recovers the published shape on the
scene it was tested on. The densification observation stays true and turns out to be insufficient, so
what it implies needs keeping straight.

\keybox{the repository's densification}{%
\texttt{densify\_and\_prune} in this codebase splits on a quantile over
accumulated \emph{absolute} gradients. The quantity it thresholds is the
accumulated magnitude of the screen-space position gradient, and the threshold
is $Q = \mathrm{quantile}(\texttt{grads\_abs},\,1-\varrho)$, where $\varrho$ is
the fraction of primitives passing the ordinary gradient test. This is the Gaussian Opacity Fields rule, and it entered the public
Mip-Splatting repository about a year \emph{after} the paper's numbers were
produced. Original 3DGS splits on $\lVert\nabla\rVert \geq \tau_{\mathrm{pos}}$
alone. So the code every reproduction of Mip-Splatting now starts from does not
densify the way the published numbers did --- a real difference from both
published methods, and one both arms of this project share.}

Two independent reasons stop it from explaining the Mip-Splatting arm's distance
from its published figure. The first is that there is no longer much
distance to explain: on the eight scenes the published mean covers, the Mip-Splatting arm is
within $\stmtArmAMaxDelta$\,dB at every test scale. The $0.50$\,dB this section
was written to account for was the artefact of \S\ref{sec:table1} --- a
seven-scene mean standing against an eight-scene published one.

The second is that the hypothesis has now been tested directly, and it fails.

\input{generated/table-pregof}

The probe reverts exactly that one change and nothing else, $45$ lines across
two files, which puts the code back in the state Mip-Splatting's published table
was produced from. The Mip-Splatting arm was re-run on it, all
$\pregofScenes$ scenes, same data, same
seed, same iterations.

\keybox{the densification explanation}{%
At the pre-GOF commit the Mip-Splatting arm reproduces the published figures to
$\pregofMaxDelta$\,dB at every test scale --- effectively exactly, and well
inside L1. If GOF densification were the cause of a standing offset, the
post-GOF column would sit clearly above this one. It does not: the effect of
that change is at most $\gofMaxEffect$\,dB and it is not consistently signed,
$+0.07$ at full resolution and $-0.11$ at $\tfrac18$.

So the densification difference is real, and it is shared by both arms, but it
is \emph{not} why the Mip-Splatting arm sits where it sits. The explanation had never been
tested, and when it finally was, it turned out to be wrong. The sentence it
replaces read: ``a better densification plausibly accounts for standing above
the published level''. That sentence was doing the work of a measurement while
being an assertion, and the word \emph{plausibly} was carrying
all of it.}

The earlier version of this section also argued the point from the 3DGS arm's
$\cOneBaseFull$\,dB on \texttt{lego} at full scale, set against a published
eight-scene mean of $33.33$. That comparison is void: \texttt{lego} is one scene
and an easy one, and a single scene standing above an eight-scene mean says
nothing about either. It is the same mistake as the table above it, at a
different denominator, and finding one is what exposed the other.

Three consequences, revised:

\begin{enumerate}[leftmargin=1.4em,itemsep=3pt]
  \item \textbf{The Mip-Splatting arm carries an L1 claim on PSNR; the 3DGS arm carries none.} The Mip-Splatting arm is
        within $\stmtArmAMaxDelta$\,dB over the eight scenes the published mean
        covers. The 3DGS arm is missing \texttt{ship}, so its column is a
        $\stmtArmBScenes$-scene mean against an eight-scene published figure and
        is not a reproduction test at all. Only PSNR is compared: the published
        SSIM and LPIPS are not tabulated per scale in the source, so the L1
        claim is scoped to PSNR and stated that way.
  \item \textbf{The between-arm comparison is now what it was claimed to be.}
        Before the opacity-compensation fix the arms differed by the 2D compensation as
        well as by the 3D filter, so the gap column mixed the method with a
        residual anti-aliasing term. Both arms now share the densification, the
        dataloader, the metrics, the LPIPS backbone \emph{and} the screen-space
        filter semantics appropriate to each, and the difference between them is
        the band-limit.
  \item \textbf{The published $\tfrac18$ gap is reproducible after all.}
        Section~\ref{sec:g2} scores it. The earlier claim that it could not be
        reproduced in this codebase was a consequence of the defect, not of the
        codebase, and it is the clearest illustration in this project of why an
        explanation that makes a failure unfixable should be distrusted until it
        has been given a chance to fail.
\end{enumerate}

The shared-densification choice stands, and now for a better reason than before.
It was defended on the grounds that pinning the pre-GOF commit would move the
difference between the arms out of the filter and back into the plumbing, which
is still true. What has changed is that the cost of keeping GOF has now been measured instead of assumed, and it is $\gofMaxEffect$\,dB at worst.

\section{Run-to-run spread}
\label{sec:spread}
Parity needs a $\sigma$ before it is a testable claim. Repeats of an identical
configuration bound the harness's own nondeterminism, and repeats at different
seeds bound the variation the initial point cloud introduces.

\input{generated/table-seedspread}

Two different quantities are at stake here, and separating them turns out to
matter.

\begin{enumerate}[leftmargin=1.4em,itemsep=3pt]
  \item \textbf{Nondeterminism at a fixed seed} --- the same configuration run
        twice. CUDA atomics in the rasteriser accumulate in a non-deterministic
        order, so bit-identical inputs do not give bit-identical outputs. The
        largest per-scale difference observed is $\repeatSpread$\,dB.
  \item \textbf{Seed-to-seed variation} --- a different initial point cloud and a
        different view order. Three seeds on two scenes, both arms, all four
        scales, give a largest sample standard deviation of
        $\mathbf{\seedSigma}$\,dB, at \seedSigmaWhere.
\end{enumerate}

The second is $\seedVsRepeat\times$ the first, so the initialisation draw matters materially more than the atomics, which is
the reason the protocol marks R5 mandatory.

\keybox{the seed spread against the L1 tolerance}{%
L1 allows $\pm0.20$\,dB. A worst-case seed standard deviation of
$\seedSigma$\,dB is about half of that, so a single-seed L1 claim at an
unfavourable scale has a margin of roughly $2\sigma$ and no more. A three-seed mean has a standard error of $\seedSigma/\sqrt{3}$ and is
comfortable; a single run is adequate, with little room to spare. Both arms' reproductions in
\S\ref{sec:table1} are single-seed results, and several of their margins are of
the same order as this $\sigma$. They should be read as L1 attained, with the margin that implies and no more.

This figure is itself a consequence of \S\ref{sec:g2}: measured on the
compensated arm it was $0.139$\,dB, and on the vanilla one it is
$\seedSigmaArmB$. The seed spread of an arm is a property of that arm.}

The spread also differs by arm: the worst case over all four scales is
$\seedSigmaArmB$\,dB for 3DGS and $\seedSigmaArmA$\,dB for Mip-Splatting. The
arm without a band-limit is the less stable one under reseeding, which is
consistent with the mechanism this thesis proposes, though two scenes are far
too few to claim it as a finding.

\section{Real scenes}
\label{sec:real}

\input{generated/table-real}

Two results here, neither of which is a PSNR.

\paragraph{The compute envelope, resolved.} Section~\ref{ch:design} recorded 16\,GB
against a stated requirement of 24\,GB, and \S8.1 of the protocol listed
``whether \texttt{counter} and \texttt{drjohnson} fit 16\,GB'' as one of only two
facts that could not be settled without running. They can be settled now:
\texttt{counter} fits, at $14\,567$\,MB. Every scene measured peaks between
$14.5$ and $\realPeakVram$\,MB against a $15\,360$\,MB card --- under a gigabyte
of headroom, at up to $\realMaxGauss$ primitives. The envelope binds, and it holds: every scene ran at full settings, with nothing
reduced to make one fit.

\paragraph{The arms agree, which is the prediction.} On every one of these
scenes the two arms land within $0.2$\,dB of each other, and they do so in both
directions: 3DGS is ahead on \texttt{bonsai}, \texttt{counter} and
\texttt{kitchen}, Mip-Splatting on \texttt{room}. That is what Proposition~\ref{prop:reduction}
plus the floor-occupancy measurement of \S\ref{sec:instrumentresults} together require: these
are well-conditioned full-orbit captures, the Nyquist floor binds, and where it
binds the two filters are identical. It is a second confirmation of the instrument result, on real data, at the level
of rendered quality, and it is why no gain is available here for a method that
improves the filter only where the floor stops binding.

\paragraph{On comparison to published figures.} None is made. The published
Mip-NeRF 360 average is over nine scenes including the five outdoor ones, which
are excluded here by the memory envelope; a four-indoor subset average is a
different quantity and setting the two side by side is precisely the error
\S\ref{ch:discussion} records as having already occurred once in this project's
earlier write-up. These rows are L2 and are the self-run control for later work.

\section{Training cost and model size}
\label{sec:measuredcost}
The estimates in the protocol's \S7 were replaced by measurement at R0, on the
same free-tier hardware every later rung uses: $\rZeroItersPerSec$ iterations per
second on a Tesla T4, $\rZeroRenderFps$ rendered test views per second, and a
peak of $\rZeroPeakVram$\,MB of VRAM. The last of these matters for scope: 16\,GB
is nowhere near binding on the Blender scenes, so the compute envelope of
\S\ref{ch:design} constrains the \emph{real} scenes.

\section{Render speed}
\label{sec:fps}

Section~\ref{sec:implcost} argues that the anisotropic filter cannot cost anything at render time:
$\Lambda_k$ is recomputed on the filter's own training schedule and never inside
the render loop. That is an argument about code, and the first question a
graphics audience asks deserves a measurement.

\input{generated/table-fps}

The anisotropic filter renders at $\bOneFpsRatio\times$ Mip-Splatting --- $\bOneFpsSlowerPct\,\%$
slower, on a model with $\bOnePrimFewerPct\,\%$ fewer primitives, which is
within the variation two training runs of the same configuration produce in
primitive count alone. The filter's own construction costs $\bOneFilterSetup$\,s,
paid once before the loop begins, and $\Lambda_k$
is never touched inside it. So the claim of \S\ref{sec:implcost} rests on a measurement.

\paragraph{The missing 3DGS row.} The benchmark mounts each arm's own
trained model and scores whichever it finds. The vanilla 3DGS model was to come from the C2 session, which ended in error with six of eight scenes complete
(\S\ref{sec:budget}); the mount did not carry a \texttt{lego} model the
benchmark could discover, so that row is absent, and it has not been estimated. Nothing
about the comparison this section defends depends on it: the quantity at issue is
the anisotropic filter against \emph{its own baseline}, and both of those are here, measured in one
session on one GPU.

\keybox{what the \texttt{render\_fps} column measures}{%
It is views divided by the wall time of a whole rendering pass,
so it carries the process start, the scene load, the PLY parse, PNG encoding and
the disk writes, which is why it reads $8$--$10$ against a published
$180$--$300$, and why quoting it as a rendering result would be a serious
misreading of this project's own data. It is retained in Table~\ref{tab:cost}
because it is what that rung logged and deleting a logged number is not an
option; it is not used as a rendering claim anywhere. The measurement that
answers the question is Table~\ref{tab:fps}, and what it defends is the
\emph{ratio}, since every cost shared by the three arms has been excluded from
all of them equally.}

\section{Floor occupancy and conditioning}
\label{sec:instrumentresults}

Both instruments are functions of the primitive positions and the capture
geometry alone, so neither needs the images, the rasteriser, or a training run
of its own. Both were therefore run over every trained scene, and over the five
capture protocols of \S\ref{sec:stress}, at a total cost of two minutes.

\input{generated/table-instruments}

\begin{figure}[htbp]
  \centering
  \includegraphics[width=0.86\linewidth]{figures/fig8-floor-binding.pdf}
  \caption{The mechanism, measured. The anisotropic filter can only differ from Mip-Splatting where
  the estimation term $s^2/\lambda_i$ exceeds the Nyquist floor $f_k^2$; below
  the dashed line Proposition~\ref{prop:reduction} makes them the same filter to
  numerical precision. On the standard full orbit the floor dominates by a factor
  of three, and it still dominates under grazing and mixed-focal capture. Only
  the arc and the cone cross the line, which is the entire content of the
  claim that this method is a capture-conditioning method. The percentages under each protocol are the fraction of
  primitives above the floor in the worst-constrained direction and in all
  directions.}
  \label{fig:floorbinding}
\end{figure}

\paragraph{Two defects sit behind these numbers.} The instruments and the
training kernels were, for a period, cloning different branches, so a protocol
named \texttt{cone} meant one capture to the instrument and another to the run
that rendered it; and the constant $0.2$ of Equation~\eqref{eq:mipfilter} was carried
on one side of a comparison and not the other, which made the estimation term
appear five times larger than it is. Both are reported in
Appendix~\ref{app:defects}; the numbers in this section are from after both were
fixed, and the agreement between the instrument and the filter's own diagnostic
is the check that confirms it.
\subsection*{The floor-occupancy instrument}

The estimation term $0.2/\lambda_i$ and Mip-Splatting's floor $f_k^2$ are the two
sides of Proposition~\ref{prop:reduction}'s hypothesis, and the ratio between
them decides whether the method can act at all. Measured across the six
protocols, it rises monotonically as the capture closes: $0.16$ on the full
orbit, $0.21$ under mixed focal, $0.24$ under grazing, $0.48$ on the arc, and
past unity only on the cone ($1.35$) and the pencil ($5.52$).

On the standard benchmark the floor dominates by a factor of six, and the
estimation term exceeds it on \textbf{zero} primitives, in \emph{any} direction,
of \emph{every} scene. By Proposition~\ref{prop:reduction} that makes the anisotropic filter \emph{be} Mip-Splatting on this benchmark, exactly:

\keybox{why the standard benchmark offers no gain}{%
The eight Blender scenes are 100-camera full orbits, the best-conditioned capture
in the suite. A method evaluated only on this benchmark could not show an
improvement, and if one appeared it would be a bug. This is the cheapest
falsification the project had, it was run before a line of method code existed,
and Table~\ref{tab:method} is the confirmation at the level of rendered quality:
$\methodParityMax$\,dB over $\methodParityScenes$ scenes and four test scales.}

\subsection*{Agreement between the two diagnostics}

The floor-occupancy instrument is computed off a trained point cloud by a NumPy instrument; the filter
prints its own occupancy during training, from the model it is actually
filtering. They are independent implementations of the same comparison, and
Table~\ref{tab:stress}'s \emph{above floor} column can be read against
Table~\ref{tab:instruments}'s \emph{worst dir} column directly.

They agree on five protocols of the six. Both read $0\,\%$ on the full orbit,
on grazing, on mixed focal and on the arc, and both put the pencil above the
floor on essentially everything, at $100\,\%$ against the filter's $93\,\%$.

The one disagreement is the cone, where the instrument reads $100\,\%$ and the
filter $0\,\%$. That looks alarming until you notice that the ratio of
estimation term to floor is $1.35$ there, which is within $35\,\%$ of unity. A
verdict balanced that finely ought to be sensitive to what it is computed on, and the two are not computed on the same thing. The instrument
takes a model trained on the \emph{full orbit} and asks what the cone's ten
cameras would have been able to determine about it. The filter takes the model
those ten cameras actually produced, which has different primitives sitting at
different depths and therefore a different floor. Each is right about its own
question. The one that decides what the filter does at training time is the filter's, so that is the one the stress table reports.

\keybox{what the constant $0.2$ bought}{%
Before the constant was corrected (\S\ref{sec:m2defect}) these two columns
disagreed on the arc as well as the cone, and by a margin no point-cloud
difference could explain. Two implementations of one comparison agreeing on five
of six protocols, and differing only where the quantity sits within a third of
its decision threshold, is a much stronger statement than either makes alone.}

\subsection*{The conditioning instrument across the stress suite}

The refutation condition at the floor-binding gate is that the floor exceeds the estimation
term on nearly every primitive in \emph{every} scene, which would make Proposition~\ref{prop:reduction} apply universally and collapse
the method to its baseline everywhere. The measurement rules that out. Degrading
the capture moves the balance as Equation~\eqref{eq:cap} requires, and the measured anisotropy tracks the
prediction across a factor of eight:

\begin{itemize}[leftmargin=1.4em,itemsep=2pt]
  \item \textbf{Full orbit} ($179^\circ$): $\instFullAniso\times$ measured
        against $\instFullCap\times$ predicted.
  \item \textbf{Grazing} ($168^\circ$): $\instGrazingAniso\times$ against
        $\instGrazingCap\times$ --- the largest relative residual in the table,
        and the protocol built to maximise obliquity, which is what
        \S\ref{sec:offaxis} predicts a parallax-only model will miss.
  \item \textbf{One-sided arc} ($94^\circ$): $\instArcAniso\times$ against
        $\instArcCap\times$.
  \item \textbf{Low-parallax cone} ($51^\circ$): $4.01\times$ against
        $3.16\times$, and the first protocol where the estimation term exceeds
        the floor along the worst-constrained direction.
  \item \textbf{Narrow pencil} ($20^\circ$): $9.42\times$ against $8.22\times$,
        and the estimation term exceeding the floor by a median factor of five.
\end{itemize}

\keybox{the cap form as the prediction}{%
Every measured anisotropy comes out above the two-view prediction of
Equation~\eqref{eq:anisotropy}, and by roughly the same factor each time. The
fact that the factor is roughly constant is the clue: this is one discrepancy observed five times, and \S\ref{sec:capprediction} says
what it is. A capture protocol is a whole distribution of cameras spread over an angular window, where the two-view form places them at its two ends. For such a distribution the
correct $\lambda_3/\lambda_1$ is asymptotically half the two-view value, which
means the two-view form understates $\sigma_D/\sigma_L$ by a factor of
$\sqrt{2}$.

Scored against the cap form instead, at each protocol's own measured span, the
$\capAgreeN$ protocols that form is built for agree to
$\capAgreeMin$--$\capAgreeMax\,\%$. The two-view expression, scored the same way
on the same measurements, falls short by as much as $\twoViewLowMax\,\%$.

The sixth protocol, \emph{grazing}, sits at $\capAgreeGrazing\,\%$, which is the expected result. The cap form
describes parallax and nothing else, and grazing is the protocol built
specifically to test obliquity at close to full parallax. It is the row the model
does not claim to cover.

One detail supports the same reading. The residual is one-signed: measurement
sits above prediction every time, never below. That is what
\S\ref{sec:offaxis} would lead you to expect, because perspective contributes
anisotropy of its own that a cap of ray directions does not contain.}

So the floor-binding gate and the anisotropy gate both pass, and the regime
table of \S\ref{sec:anisotropy} rests on measurement --- before a line of method
code exists. This is the cheapest
falsification the project had available, and it did not fire.

\paragraph{The mixed-focal prediction.} Section~\ref{sec:defects} argues that
mixing cameras is the third defect of the scalar filter, and \S\ref{sec:stress}
expected the mixed-focal protocol to expose it. The measurement shows otherwise. The anisotropy moves
from $\instFullAniso\times$ to $\instMixedAniso\times$, which is nothing, and the
median ratio of estimation term to floor from $0.16$ to $0.21$ --- a real shift,
and still a factor of five below unity. Section~\ref{sec:table3} confirms it at
the level of rendered quality on a protocol that genuinely downsamples half the
cameras: the floor binds on every primitive and the two arms differ by
$+0.017$\,dB. The honest reading is that the mixed-camera defect is a statement about the \emph{value} of $f_k$, and that a $2$--$4\times$ focal spread stays too narrow to make it bind. Either the
protocol needs a wider spread or the defect is real but small; this measurement
does not distinguish those, and it is not reported as though it did.

\section{The standard benchmark}
\label{sec:table2m}

\input{generated/table-method}

\begin{figure}[htbp]
  \centering
  \includegraphics[width=\linewidth]{figures/fig9-parity.pdf}
  \caption{Proposition~\ref{prop:reduction} as a measurement. \emph{Left:} the
  three methods over the scenes measured under all of them --- the anisotropic filter's curve is
  drawn dashed \emph{on top of} Mip-Splatting's, and is invisible beneath it
  because that is what the proposition requires. \emph{Right:} the same thing as
  a difference, against the harness's own run-to-run spread. Every scale sits
  inside $\pm\sigma$. The 3DGS curve is there to set the scale: the difference
  the band-limit makes is several decibels, and the difference between
  \emph{which} band-limit is hundredths of one.}
  \label{fig:parity}
\end{figure}

\keybox{parity, predicted and recorded first}{%
The floor-occupancy instrument measured the Nyquist floor binding on $100\,\%$ of primitives on these
captures (\S\ref{sec:instrumentresults}), and Proposition~\ref{prop:reduction}
then makes the anisotropic filter and Mip-Splatting \emph{identical}. So this table is where the method has to match its baseline exactly. A gain here
would indicate a bug, and the exact-reduction unit test is run
inside the same session that produces these numbers, aborting the run if it
fails.}

\subsection*{What parity costs}
\label{sec:methodcost}

Parity is only half a result. A method that matches its baseline exactly and
costs nothing to run is free; one that matches it exactly and costs three times
the memory has to justify itself somewhere else. Until the paired session of
\S\ref{sec:table2m} there was no run in which both methods trained the same
scenes on the same card in the same session, so cost could only have been
compared across runs --- the comparison this project has had to correct three
times.

\input{generated/table-method-cost}

\keybox{matched primitive count, higher memory}{%
The anisotropic filter densifies to within $0.2\,\%$ of Mip-Splatting's primitive count, which is
Proposition~\ref{prop:reduction} arriving from a third direction: where the
floor dominates the filter is the same filter, so the optimisation follows the
same trajectory and produces the same model. The model on disk is the same size,
and rendering is unaffected --- the filter is applied once at construction (\S\ref{sec:fps}).

Training is another matter. The anisotropic filter costs $2.8\times$ the peak VRAM and $1.32\times$
the wall-clock time, for a result that is by construction identical. The
accumulation of $\Lambda_k = \sum_n w_{nk}(J_nW_n)^{\!\top}\Sigma_{\mathrm{pix}}^{-1}(J_nW_n)$
is a $3\times3$ symmetric matrix per primitive over every training camera, and
the peak does not track the primitive count. \texttt{ficus} at $187$k primitives
and \texttt{mic} at $406$k both peak near $13$\,GB, which is what a large transient in that accumulation looks like.

That is the cost line. A practitioner reading
this table on a well-conditioned capture should not use the anisotropic filter: it is the same
model, slower, in more memory. The case for it is confined to the regime of
\S\ref{sec:table3}, and \S\ref{ch:discussion} does not argue otherwise.}

\keybox{the missing scene}{%
\texttt{ship} was scheduled in this session and only Mip-Splatting completed it.
The cause was a scheduling race: both jobs started on \texttt{ship}
simultaneously on the two GPUs, the first pair in the queue to share a scene,
and raced on the point cloud the loader writes into the source directory --- the anisotropic filter
aborted in \texttt{create\_from\_pcd} with a null point cloud, four seconds in.
That is a harness defect, since fixed by ordering the job list so that adjacent
jobs never share a scene.

It is not re-run, and the reason is this table.
Mip-Splatting alone peaks at $14\,459$\,MB on \texttt{ship}, against a
$15\,360$\,MB card. At the $2.8\times$ measured here, the anisotropic filter on \texttt{ship} does
not fit, and no ordering of the queue changes that. The paired parity figure is
therefore over $\methodParityScenes$ scenes and says so.}

\paragraph{Three bugs on the way here.} Requiring parity \emph{exactly} turned the reduction theorem into a test, and the test
failed three times before it passed --- a clamp that was doing real work, a
determinant computed in a way that lost precision at realistic scales, and an
isotropic branch that re-derived a quantity it should have taken. Each produced
plausible numbers and the wrong shape. The diagnosis of each, and the unit test
that now covers the range where they hid, is Appendix~\ref{app:defects}.
\section{The stress suite}
\label{sec:table3}

Where the claim lives. Training cameras are restricted to each protocol's subset
while the test set is left whole, so the only thing that varies between rows is
the geometry the reconstruction had available. Both methods are trained inside
one session on the same data with the same seed, so the delta is paired within scene and test scale, never taken between two
means.

\input{generated/table-stress}

\keybox{how to read this table}{%
Proposition~\ref{prop:reduction} is a constraint on this table, which is a stronger thing than a prediction about
it. Where the estimation term does not exceed the Nyquist floor in
any direction, $\Sigma_{\mathrm{filt}} = f_k^2 I$ identically and the anisotropic filter \emph{is} Mip-Splatting: the same filter, taking the same code path.
On such a row $\Delta$ is a measurement of CUDA atomics and nothing else, and a non-zero value there would be evidence of a bug. So the
diagnostic column decides whether the PSNR column is an experiment at all, which
is why it is printed by the filter during training and carried here beside the result instead of reconstructed afterwards.}

\subsection*{What the table says}

Read down the \emph{above floor} column and the $\Delta$ column together and the
result is a single sentence: \textbf{the method does nothing on five protocols
and helps on the sixth, and the five where it does nothing are the five where Proposition~\ref{prop:reduction} says it must.}

On the full orbit, the grazing orbit, the mixed-focal capture, the one-sided arc
and the count-based cone, the estimation term exceeds the Nyquist floor on
\emph{zero} primitives. $\Sigma_{\mathrm{filt}} = f_k^2 I$ identically, the mean
anisotropy is $1.00\times$ to two decimals, and the two arms take the same code
path. The measured differences over those $\stressFlatCount$ rows run from
$\stressFlatMin$ to $\stressFlatMax$\,dB --- every one inside the harness's own
run-to-run spread, and none
exceeding $0.8\sigma$. That is the proposition holding, on $\stressFlatCount$ independently degraded
captures, to the precision of the instrument.

The narrow pencil is the first capture in this project on which the filter
actually engages. The estimation term exceeds the floor on $93\,\%$ of
primitives, and the filter comes out anisotropic, at $1.47\times$.
There, and only there, the anisotropic filter is ahead. The margin is $\stressBindDelta$\,dB on the
mean, which is $\stressBindSigma$ times the run-to-run spread. That mean is made of
eight paired cells, two scenes by four test
scales, and the anisotropic filter leads in every one of them, with per-cell gains running from
$+0.163$ to $+0.217$\,dB. A mean can come out positive by luck. A sign that
holds across all eight pairings cannot.

\begin{figure}[htbp]
  \centering
  \includegraphics[width=0.92\linewidth]{figures/fig13-qualitative-pencil.pdf}
  \caption{The operating point where the anisotropic filter acts, rendered. Three training cameras
  spanning $29^\circ$; the test set is the full orbit, so most of what is
  evaluated was never observed. Numbers are PSNR for the images shown. The anisotropic filter leads on both scenes, and on all four test scales of
  each, and the lead is too small to see.}
  \label{fig:qualpencil}
\end{figure}

\keybox{how to read these renders}{%
Both reconstructions are ruined. The chair is a smear with floaters where the
unobserved sides should be, the \texttt{lego} bulldozer is barely identifiable,
and nothing in either image would be accepted as output. The anisotropic filter's advantage is real,
signed consistently, and $\stressBindSigma\times$ the run-to-run spread --- and
it is also a fifth of a decibel on a $14$\,dB reconstruction, which is invisible.

Both statements are true, and the figure is here so that the first is not read
without the second. What Table~\ref{tab:stress} establishes is that the
\emph{mechanism} is the one the derivation names: the gain appears exactly where
the estimation term overtakes the floor and nowhere else. What it does not
establish, and what this figure makes plain, is that the method is useful at
this operating point. A capture this degraded needs more cameras.}

\keybox{the predicted ordering}{%
Section~\ref{sec:anisotropy} predicts \emph{where} this method helps: it stays
inert wherever the Nyquist floor dominates, and acts more strongly as the
capture's angular span closes. A method that helped uniformly would contradict
its own derivation, and \S\ref{ch:discussion} recorded that in advance as a reason for suspicion. What the table shows is the
predicted step: five rows inside the noise, one at four times it, and the step
falls exactly at the span where the estimation term overtakes the floor. The conditioning column rises monotonically alongside it, through $1.17$,
$1.22$, $1.48$, $2.47$, $4.11$ and $6.17\times$, and is measured on the same captures the PSNR is, at the cost
\S\ref{sec:branchdefect} records. The
crossover is between the count-based cone at $54^\circ$ and the pencil at
$29^\circ$, and the suite contains both only because the two protocol
definitions of \S\ref{sec:branchdefect} disagreed and both were kept.}

\subsection*{The limits of this result}

\begin{enumerate}[leftmargin=1.4em,itemsep=3pt]
  \item \textbf{The regime is narrow, and the absolute numbers say so.} A
        three-camera capture reconstructs badly: both arms sit near $14$\,dB on
        the pencil, against $34$ on the full orbit. The gain is real and it is
        measured on a reconstruction nobody would ship. What this establishes is
        that the mechanism is the one claimed, namely that the filter acts when and only when
        the estimation term binds. It does not establish that the method is
        useful at this operating point.
  \item \textbf{Two scenes, one seed.} The eight consistent pairings are eight
        cells of one run, not eight independent experiments. The seed spread of
        \S\ref{sec:spread} bounds the noise on a cell; it does not bound the
        scene-to-scene variation of the effect, and two scenes cannot.
  \item \textbf{The count-based cone is a null result, not a missing one.} At
        $54^\circ$ the floor still binds everywhere and the method is provably
        inert. That row is evidence about where the crossover is, and it is the
        reason the crossover can be stated as a measurement instead of an extrapolation from
        the pencil alone.
\end{enumerate}

\begin{figure}[htbp]
  \centering
  \includegraphics[width=0.9\linewidth]{figures/fig7-stress-delta.pdf}
  \caption{The central empirical claim, in one panel. Paired PSNR difference
  between the anisotropic filter and Mip-Splatting per capture protocol, with the harness's own
  run-to-run spread shaded at $\pm3\sigma$; a bar inside the shaded band is not a
  result. Protocols are ordered by conditioning, best at the top, and each
  carries the fraction of primitives on which the estimation term exceeded the Nyquist floor, the quantity that determines whether the two arms
  were running the same filter at all.}
  \label{fig:stress}
\end{figure}

\section{The angular mask}
\label{sec:btwo}
\label{sec:table4}

Equation~\eqref{eq:gram} depends on primitive positions and training view
directions alone, so the criterion is measurable from a trained model with no
retraining and no GPU. It was evaluated on all eight Blender scenes and then, for
the same reason the anisotropic filter had to be, across all six capture protocols.

\keybox{the angular mask on the standard benchmark}{%
Every primitive in all eight scenes is seen from $100$ well-spread directions,
and the criterion keeps degree 3 everywhere: mean $\ell_{\max} = 3.00$,
$100\,\%$ of non-DC coefficients retained. That is the correct answer, and it is
the same shape as the floor-occupancy instrument's result for the anisotropic filter --- on a well-conditioned capture the
proposed criterion reduces to the baseline. Neither contribution can be
demonstrated on a full orbit, and this has now been measured for both, where before it was argued for one.}

\input{generated/table-b2-protocols}

\begin{figure}[htbp]
  \centering
  \includegraphics[width=0.82\linewidth]{figures/fig10-b2-retention.pdf}
  \caption{The angular mask across the capture protocols. The same shape as
  Figure~\ref{fig:floorbinding}, in the angular domain instead of the spatial one: on a well-conditioned capture the criterion keeps every coefficient and
  the method reduces to its baseline, and it removes coefficients only where the
  views cannot support them. A criterion that pruned on the full orbit would be wrong.}
  \label{fig:b2retention}
\end{figure}

Under degraded capture the criterion does what \S\ref{ch:method2} claims: a
primitive seen from a narrow cone cannot support degree 3, and is not given it.
The retained fraction falls monotonically with the conditioning, from $100\,\%$
on the full orbit to $16.6\,\%$ in the cone.

\subsection*{An amendment to Equation~\eqref{eq:lmax}}
\label{sec:b2amend}

As \S\ref{ch:method2} writes it, the criterion thresholds
$\lambda_{\min}(A_k(\ell))$ against an absolute $\tau$. Measured, that does not
work across regimes. $\lambda_{\min}$ carries units and scales with the angular
spread of the views, so one $\tau$ cannot mean the same thing on a 100-camera
orbit and a 10-camera cone. On \texttt{lego}:

\begin{table}[htbp]
  \centering
  \caption{The absolute criterion has no $\tau$ that responds gradually. Fraction
  of non-DC coefficients retained.}
  \label{tab:b2abs}
  \small
  \begin{tabular}{lrrr}
    \toprule
    protocol & $\lambda_{\min} > 0.01$ & $\lambda_{\min} > 0.001$ & $\lambda_{\min}/\lambda_{\max} > 0.01$ \\
    \midrule
    full orbit        & $100.0\,\%$ & $100.0\,\%$ & $100.0\,\%$ \\
    grazing           & $94.3\,\%$  & $100.0\,\%$ & $100.0\,\%$ \\
    one-sided arc     & $0.0\,\%$   & $82.4\,\%$  & $44.6\,\%$ \\
    low-parallax cone & $0.0\,\%$   & $20.0\,\%$  & $13.3\,\%$ \\
    \bottomrule
  \end{tabular}
\end{table}

At $\tau = 0.01$ the absolute form keeps everything on a full orbit and
\emph{nothing} on either degraded protocol; at $\tau = 0.001$ it keeps $82\,\%$
on the arc, which is nearly everything. It is a cliff.
The condition number $\lambda_{\min}/\lambda_{\max}$ is scale-free and behaves as
the criterion should at a single threshold, which is why it is the default in the
implementation and why Equation~\eqref{eq:lmax} should be restated in that form.
This is a change to the method that only measurement could have produced, and it is recorded as such instead of folded silently into the derivation.

\subsection*{The full-orbit comparison}

\input{generated/table-b2}

Every cell in that table is identical across the three rows, and that is the result. On the standard benchmark the
criterion retains $100\,\%$ of non-DC coefficients, so the masked model \emph{is} the control, the same file scored through the same
render and metrics path, and magnitude pruning matched to the same fraction
prunes nothing either. The $\tau$ sweep says the same thing from another direction:
between $\tau = 0.003$ and $\tau = 0.1$, a factor of thirty, nothing on a full
orbit changes.

\keybox{the comparison this table cannot make}{%
Table~\ref{tab:b2} exists to ask whether identifiability masking beats magnitude
pruning \emph{at matched model size}. That question has no content when the
criterion prunes nothing: both rows keep everything, so both score identically,
and the table cannot distinguish a good criterion from a bad one. Reporting the
three equal rows as evidence that the angular mask matches magnitude pruning would be the error this chapter warns about elsewhere --- a condition tested where it was
constrained to be satisfied. So the comparison was run again on a capture where
the criterion does mask.}

\subsection*{The comparison on a capture where the angular mask prunes}
\label{sec:b2loses}

Restricting the cameras the \emph{criterion} reads to the low-parallax cone, and
leaving the test set whole, the angular mask retains $13$--$20\,\%$ of non-DC coefficients.
Magnitude pruning is then matched to that exact fraction, so the two models are
the same size and differ only in which coefficients survive.

\input{generated/table-b2-matched}

\keybox{the angular mask against magnitude pruning}{%
On every configuration where both methods actually prune, magnitude pruning
wins. Against the free control it wins by $0.60$ and $0.82$\,dB on \texttt{lego}
and by $1.82$\,dB on \texttt{chair}; against the \emph{fair} control below, by
$0.21$ and $0.42$\,dB on \texttt{lego} and by $0.09$\,dB on \texttt{chair}.
Chapter~\ref{ch:method2} states the standard this was to be held to: ``a criterion that cannot beat
magnitude pruning at matched size is not worth the machinery''. By that standard
the criterion fails, on the fair comparison
as well as the unfair one. The row at $0\,\%$ is both methods discarding
everything and agreeing by definition; it is not evidence either way.}

\paragraph{Why, and what survives.} Identifiability and importance are different
quantities, and the metric rewards the second. A coefficient can be badly
conditioned by the capture and still carry most of the fitted energy for that
primitive; zeroing it costs more PSNR than zeroing a well-determined coefficient
that happens to be tiny. Magnitude pruning optimises directly for what PSNR measures; the angular mask
optimises for something else and is then scored on PSNR.

\paragraph{The confound in the first comparison.} A second asymmetry sits in the
free control, and it is a large one: the angular mask masks whole
\emph{degrees}, removing coefficients in blocks of $2\ell+1$ nested from the top
down, while free magnitude pruning ranks individual coefficients and keeps the
largest anywhere. At the same budget that is strictly more freedom, so a loss against it is evidence about the \emph{structure} of the mask, when
the criterion is what is on trial. The first version of this
section noted the confound and left it, which was the right thing to notice and
the wrong place to stop.

The control that separates them is magnitude pruning restricted to the angular mask's own
nested degrees: same structure, same budget, and the only remaining difference
is what the two rank by. It is the \emph{blocked} column of
Table~\ref{tab:btwomatched}, and it moves the answer a long way. On
\btwoStructureScene, where the free control had beaten the angular mask by
$\btwoFreeGap$\,dB, the fair one beats it by $\btwoBlockedGapMin$ --- inside the
seed spread of $\seedSigma$\,dB, so on that scene the criterion and magnitude pruning are indistinguishable. $\btwoStructureShare\,\%$ of what looked
like a verdict on identifiability was a verdict on block structure.

What does not change is the direction. On \texttt{lego} the fair control still
wins at both budgets where anything is pruned, by up to
$\btwoBlockedGapMax$\,dB --- four times the seed spread. So the honest statement is narrower and
firmer than the one it replaces: at matched structure and matched budget,
ranking spherical-harmonic coefficients by how well the capture determines them
is worse than ranking them by how large they are --- by a fifth of a decibel
instead of two, decisive on one scene of the two and inside the noise on the
other.

What survives is narrower than what was claimed. The criterion still does what \S\ref{sec:b2scale} says it does: it responds
gradually to the conditioning, retaining $100\,\%$, $62\,\%$ and $25\,\%$ as the
capture degrades. That is a statement about the \emph{geometry}, measured and settled. What is not
supported is the step from there to rendered quality at matched size. On the
evidence here, identifiability is a good description of what the capture
determines and a poor rule for deciding what to keep.

\paragraph{Model size on disk.} Model size holds still, and has to: masking coefficients to zero leaves a fixed-width PLY the same size on
disk. Realising the saving needs variable-length SH storage, which
Chapter~\ref{ch:plan} schedules for Spring. The $74.8$\,MB in all three rows is
the same file written three times.

\section{What is established}

The reproduction rungs establish a control. Both baselines now exist as this project's own
measurements, on known hardware, from pinned commits, with the arm separation verified on the saved models. Every later comparison is against these numbers, which is
the L3 discipline of Table~\ref{tab:levels}.

About the proposed method they establish four things and refute one. The anisotropic filter is
implemented and reduces to its baseline exactly where the theory says it must,
verified to $1.4\times10^{-21}$ in the unit test and to $\methodParityMax$\,dB, comfortably below the seed spread, over
$\methodParityScenes$ scenes and four test scales in paired training runs. The conditioning both methods depend on
varies across capture protocols as Equation~\eqref{eq:cap} predicts.
The stress suite has been run, and the prediction recorded before it holds: on
the $\stressFlatCount$ protocols where the Nyquist floor still dominates the
difference is $\stressFlatMin$ to $\stressFlatMax$\,dB and the filter is
isotropic to two decimals, and on the one where the estimation term binds the anisotropic filter
leads by $\stressBindDelta$\,dB, on the mean and on every paired cell. And the cost of that has now been measured, where before it was assumed: $2.8\times$ peak VRAM and
$1.32\times$ training time for a model with the same primitive count
(\S\ref{sec:methodcost}).

What is refuted is the angular mask. At matched structure and matched budget it loses to
magnitude pruning, by $\btwoBlockedGapMax$\,dB at worst, and the criterion is
reported as failing the standard Chapter~\ref{ch:method2} set for it
(\S\ref{sec:btwo}).

What is \emph{not} established is that the anisotropic filter is useful. The regime where it helps
is one in which both methods produce reconstructions nobody would ship
(Figure~\ref{fig:qualpencil}), and on well-conditioned captures it is the same
model for three times the memory. Chapter~\ref{ch:discussion} takes that up;
this chapter's job is the measurement, and the measurement is that the mechanism
is the one the derivation names.

\section{Compute expenditure}
\label{sec:budget}

Every number in this chapter was produced on free cloud GPUs, within a weekly
allowance that the project exhausted more than once. The ladder was designed around that constraint. The accounting of which rung
cost what, which session died and why, and what the constraint forbade outright
is set out in Appendix~\ref{app:harness}. The short version is that compute was the binding constraint throughout.